% !PRINTED 342
\weberpage{342}
%<<CONT
are representable with integral $t_1, t_2$ (Bd.\,II, §\,163, 164). For
then, if $\alpha$ is representable by $\lambda$, every product $\omega\alpha$ is
also representable by $\lambda$. In our case this condition reduces to the requirement that $\alpha_1\theta, \alpha_2\theta$ must be representable by $\lambda$. But if
$(a, b, c)$ is any quadratic form of the discriminant $\Delta$
with \gesperrt{positive first} coefficient $a$, then we may also take
$\tfrac12(b+\sqrt{\Delta})$ for $\theta$, and since
\[
\left(\frac{b+\sqrt{\Delta}}{2}\right)^2 = -ac + b\,\frac{b+\sqrt{\Delta}}{2}
\]
holds, (11) is always the basis of an ideal in $\Omega$. Here
the positive number $a_{22}$ may be assumed arbitrarily. If we set
$a_{22} = 1$, then
\weq{15}{\lambda = a(x-\omega y)}
is the basis form of an ideal $\mathfrak{a}$ completely determined by $(a, b, c)$, whose
norm is equal to $a$, and a basis of this ideal is
\weq{16}{(a,\ a\omega).}

\begin{center}
§ 92. \gesperrt{Prime Ideals in the Quadratic Field.}
\end{center}

A prime $p$ is in the quadratic field $\Omega$ either
still indecomposable itself, and is then a prime ideal of the second
degree, or $p$ splits into two prime ideals of the first degree $\mathfrak{p}, \mathfrak{p}'$.
We have accordingly two cases to distinguish:
\[
\text{a)}\quad p = \mathfrak{p} \text{ in } \Omega \text{ indecomposable:}\quad N(\mathfrak{p}) = p^2,
\]
\[
\text{b)}\quad p = \mathfrak{p}\mathfrak{p}' \qquad\qquad N(\mathfrak{p}) = N(\mathfrak{p}') = p,
\]
and in the latter case the two prime ideals $\mathfrak{p}, \mathfrak{p}'$ may be either
different from one another or else equal, so that we can distinguish
yet a third case:
\[
\text{c)}\quad p = \mathfrak{p}^2,\quad N(\mathfrak{p}) = p.
\]
It is now a question of ascertaining the conditions
under which the one or the other of these cases occurs.

If $N(\mathfrak{p}) = p$, then $p$ is the number of incongruent
numbers $(\mathrm{mod}\ \mathfrak{p})$ in $\Omega$, and consequently in this case \gesperrt{every number}
in $\Omega$ is congruent to one of the rational numbers $r$:
\weq{1}{0,\ 1,\ 2,\ \ldots,\ p-1}

If on the other hand $N(\mathfrak{p}) = p^2$, then there are also numbers which,
with respect to $\mathfrak{p}$, are not congruent to any rational number.
